% \iffalse meta-comment
%
% hit-footnote.dtx 是各个类共用的脚注层
%
% \fi
%
% \section{共享脚注层}
%
% \changes{v4.0a}{2026/10/08}{画圈的几何照 Typst \pkg{quan} 0.3.0 的 \texttt{digit-fit} 重写：圈径与圈线粗细取自宋体
%   角色字体的 \texttt{⑩}，圈里数字按分组缩放、偏移；内建宋体查表，用户指定的字体现场量；
%   数字一律用宋体角色字体（\hologo{LuaLaTeX} 下也是）；不再依赖 \pkg{circledsteps}，改用 \pkg{pict2e}}
% \changes{v4.0a}{2026/09/27}{带圈脚注号按码位排中文字走引擎层：\hologo{pdfLaTeX} 下 \cs{char} 排不了 255 以上的码位}
% \changes{v4.0a}{2026/08/15}{脚注层抽出来共用，报告的脚注由此也是带圈编号、小五宋体；
%   \cs{quan} 先取当前宋体的带圈字符，超出才自己画（几何按字体自带的
%   \texttt{U+2460} 定）；\cs{@makefntext} 改成包一层（\cs{pretocmd} 在 expl3 里会吃掉
%   字面空格，脚注整块上移 4.2bp）；\option{layout}/\option{bottom} 的两支各发
%   \cs{raggedbottom} 或 \cs{flushbottom}，推迟到 \texttt{begindocument}}
%
% 提供带圈数字及报告脚注样式，终稿脚注由 docxlike 处理。
%
% 分成两截。\cs{hit@setup@footnote@packages} 只管加载 \pkg{footmisc}，位置要留在依赖模块
% 里不能动；\cs{hit@setup@footnote} 管余下的样式，得排在 \pkg{footmisc} 与字号都就位
% 之后调用，因为它要接管 \cs{@makefntext} 并存一份当时的 \cs{footnotesize}。
%
%    \begin{macrocode}
%<*shared-footnote>
%    \end{macrocode}
%
% \begin{macro}{\hit@setup@footnote@packages}
% 终稿的脚注由 \pkg{docxlike} 照 Word 排（见 \file{hit-final-footnote.dtx}），不加载 \pkg{footmisc}；
% 下面说的是报告。脚注按页编号。页底策略分两支：\texttt{ragged} 不填补页底空白，脚注仍要压在
% 底部，所以 \pkg{footmisc} 要带 \texttt{bottom}；\texttt{flush} 把空白摊到
% 行间，脚注自然就在底部，不用带。
%
% 发命令的时机与加载宏包的时机不一样。宏包选项只能在加载时定，而
% \cs{raggedbottom} 与 \cs{flushbottom} 是随时可改的，所以推迟到
% \texttt{begindocument}——不然 \cs{hitsetup} 里写的 \option{bottom} 比这里晚，
% 改了标志位却已经来不及。
%    \begin{macrocode}
\cs_new_protected:Npn \hit@setup@footnote@packages
  {
    \bool_if:NF \g__hit_final_bool
      {
        \bool_if:NTF \g__hit_raggedbottom_bool
          { \RequirePackage [ bottom , perpage , hang ] { footmisc } }
          { \RequirePackage [ perpage , hang ] { footmisc } }
      }
    \AtBeginDocument
      {
        \bool_if:NTF \g__hit_raggedbottom_bool
          { \raggedbottom }
          { \flushbottom }
      }
  }
%    \end{macrocode}
% \end{macro}
%
% \section{带圈数字}
% \begin{macro}{\quan}
%
% \cs{quan} 先取当前宋体里现成的带圈字符，字体没有才自己画。画出来的圈要与字体自带的
% \texttt{①}--\texttt{⑩} 对得上，口径与 Typst 的 \pkg{quan} 0.3.0（\texttt{digit-fit}）一致：
% 圈径、圈线粗细取自参照字形，圈里数字的高、宽、偏移按数字分组定。
%
% 圈里的数字与参照字形一律取宋体角色（\texttt{songti}）的字体，不跟当前的中文字体走：
% \hologo{LuaLaTeX} 下 ASCII 数字被 \pkg{luatexja} 留在西文字符范围，切不到和文字体，所以
% 用 \pkg{fontspec} 按宋体角色单独声明一个字族，两个引擎行为一致。字族在第一次画圈时才声明，
% 那时字体角色表已经定稿。
%
% 数据分两级。
% \begin{itemize}
%   \item 内建字体（预设里的宋体：SimSun、Songti~SC、FZShuSong、FandolSong、Noto~Serif~CJK~SC）
%     直接查表，表是 \pkg{quan} 0.3.0 离线标定的，只收「圈码字体 = 数字字体」的那一列。
%   \item 用户自己指定的字体表里没有，现场量：圈径、数字的墨迹高宽都从实际字形取，
%     数字高宽比例用 \pkg{quan} 的通用近似（\texttt{\_FIT-APPROX}，各圈码来源等权汇总），
%     圈线粗细取 \texttt{0.0315em}（粗体乘 1.45）。圈线粗细读不到字形轮廓
%     （\pkg{luaotfload} 不保留轮廓），只能退到这个常数。
% \end{itemize}
%
% 表项为 \texttt{字体/字重/分组}，值依次是水平缩放、数字高占圈高的比例、水平偏移、垂直偏移
% （后两项以圈高为单位，垂直向下为正）、小字号（10bp 及以下）与大字号（24bp 及以上）的宽度系数。
% 分组是：单个 1、其余单个数字、11、含 1 的两位数、其余两位数。
%    \begin{macrocode}
\prop_const_from_keyval:Nn \c__hit_quan_fit_prop
  {
    fandolsong/regular/single-1 = {1.0000,0.6400,0.0000,0.0200,1.000,1.000},
    fandolsong/regular/single = {1.0000,0.6600,0.0000,0.0300,1.000,1.000},
    fandolsong/regular/all-ones = {0.7485,0.6600,0.0175,0.0300,0.688,0.872},
    fandolsong/regular/two-with-1 = {0.7485,0.6600,0.0175,0.0300,0.611,0.780},
    fandolsong/regular/two = {0.7485,0.6600,0.0175,0.0300,0.611,0.780},
    fandolsong/bold/single-1 = {1.0000,0.6300,0.0000,-0.0350,1.000,1.000},
    fandolsong/bold/single = {1.0000,0.6400,0.0000,-0.0300,1.000,1.000},
    fandolsong/bold/all-ones = {0.7485,0.6600,0.0175,0.0300,0.587,0.796},
    fandolsong/bold/two-with-1 = {0.7485,0.6600,0.0175,0.0300,0.546,0.712},
    fandolsong/bold/two = {0.7485,0.6600,0.0175,0.0300,0.546,0.712},
    fzshusong-z01s/regular/single-1 = {1.1636,0.6923,0.0000,0.0055,1.000,1.000},
    fzshusong-z01s/regular/single = {1.0430,0.7060,0.0000,0.0069,0.878,1.000},
    fzshusong-z01s/bold/single-1 = {1.1636,0.6923,0.0000,0.0055,1.000,1.000},
    fzshusong-z01s/bold/single = {1.0430,0.7060,0.0000,0.0069,0.747,1.000},
    noto-serif-cjk-sc/regular/single-1 = {0.8986,0.5674,-0.0056,-0.0140,1.000,1.000},
    noto-serif-cjk-sc/regular/single = {0.8104,0.5730,0.0028,-0.0084,1.000,1.000},
    noto-serif-cjk-sc/regular/all-ones = {0.6101,0.5674,-0.0028,-0.0140,1.000,1.000},
    noto-serif-cjk-sc/regular/two-with-1 = {0.6086,0.5758,0.0084,-0.0098,1.000,1.000},
    noto-serif-cjk-sc/regular/two = {0.6406,0.5787,-0.0042,-0.0084,0.911,0.974},
    noto-serif-cjk-sc/bold/single-1 = {0.8889,0.5611,-0.0056,-0.0139,1.000,1.000},
    noto-serif-cjk-sc/bold/single = {0.7905,0.5722,0.0000,-0.0056,1.000,1.000},
    noto-serif-cjk-sc/bold/all-ones = {0.6358,0.5611,-0.0028,-0.0139,1.000,1.000},
    noto-serif-cjk-sc/bold/two-with-1 = {0.6027,0.5722,0.0069,-0.0083,0.945,1.000},
    noto-serif-cjk-sc/bold/two = {0.6097,0.5722,-0.0042,-0.0083,0.819,0.924},
    simsun/regular/single-1 = {0.9455,0.5989,0.0000,0.0028,1.000,1.000},
    simsun/regular/single = {1.0998,0.6130,-0.0056,0.0000,1.000,1.000},
    simsun/regular/all-ones = {0.6125,0.6045,0.0141,-0.0056,1.000,1.000},
    simsun/regular/two-with-1 = {0.6125,0.6045,0.0141,-0.0056,1.000,1.000},
    simsun/regular/two = {0.6125,0.6045,0.0141,-0.0056,1.000,1.000},
    simsun/bold/single-1 = {0.9455,0.5989,0.0000,0.0028,1.000,1.000},
    simsun/bold/single = {1.0998,0.6130,-0.0056,0.0000,1.000,1.000},
    simsun/bold/all-ones = {0.6125,0.6045,0.0141,-0.0056,1.000,1.000},
    simsun/bold/two-with-1 = {0.6125,0.6045,0.0141,-0.0056,0.945,1.000},
    simsun/bold/two = {0.6125,0.6045,0.0141,-0.0056,0.945,1.000},
    songti-sc/regular/single-1 = {1.2909,0.7039,0.0028,-0.0307,1.000,1.000},
    songti-sc/regular/single = {1.0000,0.7039,0.0000,-0.0140,1.000,1.000},
    songti-sc/regular/all-ones = {0.7862,0.7062,0.0000,-0.0282,0.882,1.000},
    songti-sc/regular/two-with-1 = {0.7143,0.7401,0.0056,-0.0169,0.646,0.748},
    songti-sc/regular/two = {0.7143,0.7401,0.0056,-0.0169,0.550,0.681},
    songti-sc/bold/single-1 = {0.9167,0.7017,-0.0166,-0.0166,1.000,1.000},
    songti-sc/bold/single = {1.0310,0.7348,-0.0083,0.0000,0.867,1.000},
    songti-sc/bold/all-ones = {0.6400,0.6906,-0.0028,-0.0166,0.942,1.000},
    songti-sc/bold/two-with-1 = {0.6605,0.7072,0.0221,-0.0069,0.571,0.783},
    songti-sc/bold/two = {0.6824,0.7017,0.0193,-0.0166,0.459,0.682},
  }
%    \end{macrocode}
%
% 原生 \texttt{①}--\texttt{⑩} 的圈线厚度与外径之比，依次为常规与粗体。别名把地区变体
% 并到同一副字体上（只标定过 SC 版，其余沿用），\texttt{FZShuSong-Z01} 与标定表里的
% \texttt{fzshusong-z01s} 是同一副。键是字体名小写、空格换成连字符。
%    \begin{macrocode}
\prop_const_from_keyval:Nn \c__hit_quan_stroke_prop
  {
    simsun = {0.03530,0.03530},
    libertinus-serif = {0.05589,0.06587},
    noto-serif-cjk-sc = {0.03123,0.03553},
    noto-sans-cjk-sc = {0.03412,0.03910},
    songti-sc = {0.03421,0.06352},
    fzshusong-z01s = {0.02877,0.02877},
    fandolsong = {0.03827,0.03829},
    heiti-sc = {0.06343,0.06343},
    pingfang-sc = {0.03815,0.04839},
    microsoft-yahei = {0.06483,0.06483},
    simhei = {0.04881,0.04881},
    kaiti = {0.04881,0.04881},
    fangsong = {0.03513,0.03513},
    hiragino-sans-gb = {0.03450,0.04383},
    stsong = {0.02805,0.02805},
  }
\prop_const_from_keyval:Nn \c__hit_quan_alias_prop
  {
    nsimsun = simsun,
    fzshusong-z01 = fzshusong-z01s,
    noto-serif-cjk-tc = noto-serif-cjk-sc,
    noto-serif-cjk-jp = noto-serif-cjk-sc,
    noto-serif-cjk-kr = noto-serif-cjk-sc,
    source-han-serif-sc = noto-serif-cjk-sc,
    source-han-serif = noto-serif-cjk-sc,
    noto-sans-cjk-tc = noto-sans-cjk-sc,
    noto-sans-cjk-jp = noto-sans-cjk-sc,
    noto-sans-cjk-kr = noto-sans-cjk-sc,
    source-han-sans-sc = noto-sans-cjk-sc,
    source-han-sans = noto-sans-cjk-sc,
  }
%    \end{macrocode}
%
% 表里没有的字体用的通用近似：各圈码来源等权汇总的（水平缩放，数字高占内圈的比例）。
%    \begin{macrocode}
\prop_const_from_keyval:Nn \c__hit_quan_approx_prop
  {
    regular/single-1 = {0.9803,0.6338},
    regular/single = {0.9178,0.6471},
    regular/all-ones = {0.6018,0.6063},
    regular/two-with-1 = {0.6325,0.6341},
    regular/two = {0.6440,0.5798},
    bold/single-1 = {0.9047,0.6202},
    bold/single = {0.9158,0.6354},
    bold/all-ones = {0.5754,0.5782},
    bold/two-with-1 = {0.6078,0.6057},
    bold/two = {0.6383,0.5909},
  }
%    \end{macrocode}
%
% 变量。参照字形的墨迹高、深写在全局变量里，因为是在装圈号字体的组里量的。
%    \begin{macrocode}
\box_new:N \l__hit_quan_digits_box
\box_new:N \l__hit_quan_ring_box
\dim_new:N \g__hit_quan_ref_ht_dim
\dim_new:N \g__hit_quan_ref_dp_dim
\int_new:N \l__hit_quan_value_int
\int_new:N \l__hit_quan_slot_int
\bool_new:N \l__hit_quan_bold_bool
\bool_new:N \l__hit_quan_cal_bool
\bool_new:N \l__hit_quan_fitted_bool
\bool_new:N \l__hit_quan_first_bool
\bool_new:N \g__hit_quan_font_tried_bool
\tl_new:N \g__hit_quan_fontfam_tl
\tl_new:N \g__hit_quan_family_tl
\tl_new:N \l__hit_quan_fontfam_tl
\tl_new:N \l__hit_quan_tmpx_tl
\tl_new:N \l__hit_quan_tmpy_tl
\tl_new:N \l__hit_quan_tmpd_tl
\tl_new:N \l__hit_quan_entry_tl
\tl_new:N \l__hit_quan_group_tl
\tl_new:N \l__hit_quan_key_tl
\str_new:N \l__hit_quan_digits_str
\fp_new:N \l__hit_quan_fs_fp
\fp_new:N \l__hit_quan_h_fp
\fp_new:N \l__hit_quan_leg_fp
\fp_new:N \l__hit_quan_t_fp
\fp_new:N \l__hit_quan_w_fp
\fp_new:N \l__hit_quan_inner_fp
\fp_new:N \l__hit_quan_avail_fp
\fp_new:N \l__hit_quan_clear_fp
\fp_new:N \l__hit_quan_wd_fp
\fp_new:N \l__hit_quan_hd_fp
\fp_new:N \l__hit_quan_sx_fp
\fp_new:N \l__hit_quan_sy_fp
\fp_new:N \l__hit_quan_dx_fp
\fp_new:N \l__hit_quan_dy_fp
\fp_new:N \l__hit_quan_tmp_fp
%    \end{macrocode}
%
% \texttt{n} 对应的带圈数字码位：1--20 在 \texttt{U+2460}，21--35 在
% \texttt{U+3251}，36--50 在 \texttt{U+32B1}，超出范围返回 0。
%    \begin{macrocode}
\cs_new:Npn \__hit_quan_slot:n #1
  {
    \int_compare:nNnTF {#1} < { 1 } { 0 }
      {
        \int_compare:nNnTF {#1} < { 21 } { \int_eval:n { #1 - 1 + "2460 } }
          {
            \int_compare:nNnTF {#1} < { 36 } { \int_eval:n { #1 - 21 + "3251 } }
              {
                \int_compare:nNnTF {#1} < { 51 }
                  { \int_eval:n { #1 - 36 + "32B1 } } { 0 }
              }
          }
      }
  }
%    \end{macrocode}
%
% 把当前字体设成宋体角色的圈号字族。不能在 \pkg{pdfLaTeX} 下用（那边中文字体不走 \pkg{fontspec}），
% 字族没声明出来时什么也不做，数字就跟当前字体。
%    \begin{macrocode}
\cs_new_protected:Npn \__hit_quan_digit_font:
  {
    \tl_if_empty:NF \g__hit_quan_fontfam_tl
      { \fontfamily { \g__hit_quan_fontfam_tl } \fontshape { n } \selectfont }
  }
%    \end{macrocode}
%
% 第一次画圈时按宋体角色声明圈号字族，同时定下查表用的字体键。选项取与中文宋体相同的那一份，
% 去掉文字基线校正（那是给汉字字面用的，数字与圈是各自对齐的，不能让数字单独上下挪）。
% 再关掉 \texttt{locl}：思源宋体（Noto Serif CJK）在拉丁文字（\texttt{latn}）下用 \texttt{locl}
% 把阿拉伯数字换成另一套字形，“1”的宽从 0.471\,em 变成 0.550\,em，与 Typst 排出来的
% 默认数字不同，校准表就对不上了。
%    \begin{macrocode}
\cs_new_protected:Npn \__hit_quan_font_ensure:
  {
    \bool_if:NF \g__hit_quan_font_tried_bool
      {
        \bool_gset_true:N \g__hit_quan_font_tried_bool
        \sys_if_engine_pdftex:F
          {
            \cs_if_exist:NT \__hit_font_cjk_role:n
              {
                \group_begin:
                  \__hit_font_cjk_role:n { songti }
                  \regex_replace_once:nnN
                    { ,? RawFeature = \{ [^\}]* \} } { } \l__hit_font_role_options_tl
                  \tl_if_blank:VF \l__hit_font_role_options_tl
                    { \tl_put_right:Nn \l__hit_font_role_options_tl { , } }
                  \tl_put_right:Nn \l__hit_font_role_options_tl { RawFeature = { -locl } }
                  \use:e
                    {
                      \exp_not:N \fontspec_set_family:Nnn \exp_not:N \l__hit_quan_fontfam_tl
                        { \exp_not:V \l__hit_font_role_options_tl }
                        { \exp_not:V \l__hit_font_role_name_tl }
                    }
                  \tl_gset_eq:NN \g__hit_quan_fontfam_tl \l__hit_quan_fontfam_tl
                  \str_set:Nx \l_tmpa_str { \str_foldcase:V \l__hit_font_role_name_tl }
                  \str_replace_all:Nnn \l_tmpa_str { ~ } { - }
                  \prop_get:NVNF \c__hit_quan_alias_prop \l_tmpa_str \l_tmpb_str
                    { \str_set_eq:NN \l_tmpb_str \l_tmpa_str }
                  \tl_gset_eq:NN \g__hit_quan_family_tl \l_tmpb_str
                \group_end:
              }
          }
      }
  }
%    \end{macrocode}
%
% 量参照字形（\texttt{⑩}，没有就用 \texttt{①}）的墨迹高与深。\hologo{LuaLaTeX} 从当前字体的字符表取，
% 因为盒子的高深在和文字体那边是字面框的值，不是墨迹；\hologo{XeLaTeX} 与 \hologo{pdfLaTeX} 直接量盒子。
% 调用时当前字体必须已经是圈号字体。
%    \begin{macrocode}
\sys_if_engine_luatex:TF
  {
    \cs_new_protected:Npn \__hit_quan_ink_set:w #1 , #2 \q_stop
      {
        \dim_gset:Nn \g__hit_quan_ref_ht_dim { #1 sp }
        \dim_gset:Nn \g__hit_quan_ref_dp_dim { #2 sp }
      }
    \cs_new_protected:Npn \__hit_quan_ink:n #1
      {
        \use:e
          {
            \exp_not:N \__hit_quan_ink_set:w
            \tex_directlua:D
              {
                local~f~=~font.getfont(font.current())~
                local~c~=~f~and~f.characters~and~f.characters[\int_eval:n {#1}]~
                if~c~and~c.height~then~
                  tex.print(math.floor(c.height)..','..math.floor(math.max(c.depth~or~0,0)))
                ~else~
                  tex.print('0,0')
                ~end
              }
            \exp_not:N \q_stop
          }
      }
  }
  {
    \cs_new_protected:Npn \__hit_quan_ink:n #1
      {
        \hbox_set:Nn \l__hit_quan_ring_box { \hit@engine@char:n {#1} }
        \dim_gset:Nn \g__hit_quan_ref_ht_dim { \box_ht:N \l__hit_quan_ring_box }
        \dim_gset:Nn \g__hit_quan_ref_dp_dim { \box_dp:N \l__hit_quan_ring_box }
      }
  }
\cs_new_protected:Npn \__hit_quan_measure_ref:
  {
    \group_begin:
      \__hit_quan_digit_font:
      \__hit_quan_ink:n { "2469 }
      \dim_compare:nNnT { \g__hit_quan_ref_ht_dim + \g__hit_quan_ref_dp_dim } < { 0.1pt }
        { \__hit_quan_ink:n { "2460 } }
    \group_end:
    \dim_compare:nNnT { \g__hit_quan_ref_ht_dim + \g__hit_quan_ref_dp_dim } < { 0.1pt }
      {
        \dim_gset:Nn \g__hit_quan_ref_ht_dim { \fp_to_dim:n { 0.824 * \dim_to_fp:n { \f@size pt } } }
        \dim_gset:Nn \g__hit_quan_ref_dp_dim { \fp_to_dim:n { 0.113 * \dim_to_fp:n { \f@size pt } } }
      }
  }
%    \end{macrocode}
%
% 数字分组，与 \pkg{quan} 的校准脚本一致。
%    \begin{macrocode}
\cs_new_protected:Npn \__hit_quan_group:
  {
    \int_compare:nNnTF { \l__hit_quan_value_int } < { 10 }
      {
        \int_compare:nNnTF { \l__hit_quan_value_int } = { 1 }
          { \tl_set:Nn \l__hit_quan_group_tl { single-1 } }
          { \tl_set:Nn \l__hit_quan_group_tl { single } }
      }
      {
        \int_compare:nNnTF { \l__hit_quan_value_int } = { 11 }
          { \tl_set:Nn \l__hit_quan_group_tl { all-ones } }
          {
            \str_if_in:NnTF \l__hit_quan_digits_str { 1 }
              { \tl_set:Nn \l__hit_quan_group_tl { two-with-1 } }
              { \tl_set:Nn \l__hit_quan_group_tl { two } }
          }
      }
  }
\cs_new:Npn \__hit_quan_entry:n #1
  { \clist_item:Nn \l__hit_quan_entry_tl {#1} }
%    \end{macrocode}
%
% 通用近似：数字之间加 \texttt{0.04em} 字距，高度按内圈的比例定，宽度不超出内圈的弦；
% 三位及以上的数不分高宽，等比缩到放进内圈对角线。
%    \begin{macrocode}
\cs_new_protected:Npn \__hit_quan_fit_approx:
  {
    \hbox_set:Nn \l__hit_quan_digits_box
      {
        \__hit_quan_digit_font:
        \bool_set_true:N \l__hit_quan_first_bool
        \str_map_inline:Nn \l__hit_quan_digits_str
          {
            \bool_if:NTF \l__hit_quan_first_bool
              { \bool_set_false:N \l__hit_quan_first_bool }
              { \tex_kern:D 0.04em \scan_stop: }
            ##1
          }
      }
    \fp_set:Nn \l__hit_quan_wd_fp { \dim_to_fp:n { \box_wd:N \l__hit_quan_digits_box } }
    \fp_set:Nn \l__hit_quan_hd_fp
      { \dim_to_fp:n { \box_ht:N \l__hit_quan_digits_box + \box_dp:N \l__hit_quan_digits_box } }
    \fp_set:Nn \l__hit_quan_inner_fp { max ( 0.01 , \l__hit_quan_inner_fp ) }
    \int_compare:nNnTF { \l__hit_quan_value_int } < { 100 }
      { }
      { \tl_set:Nn \l__hit_quan_group_tl { two } }
    \tl_set:Nx \l__hit_quan_key_tl
      { \bool_if:NTF \l__hit_quan_bold_bool { bold } { regular } / \l__hit_quan_group_tl }
    \prop_get:NVNT \c__hit_quan_approx_prop \l__hit_quan_key_tl \l__hit_quan_entry_tl
      {
        \fp_set:Nn \l__hit_quan_tmp_fp
          { min ( \__hit_quan_entry:n { 2 } * \l__hit_quan_h_fp , 0.8 * \l__hit_quan_inner_fp ) }
        \fp_set:Nn \l__hit_quan_sy_fp { \l__hit_quan_tmp_fp / \l__hit_quan_hd_fp }
        \int_compare:nNnTF { \l__hit_quan_value_int } < { 100 }
          {
            \fp_set:Nn \l__hit_quan_sx_fp
              {
                min
                  (
                    \__hit_quan_entry:n { 1 } ,
                    sqrt ( max ( 0 , \l__hit_quan_inner_fp ^ 2 - \l__hit_quan_tmp_fp ^ 2 ) )
                      / \l__hit_quan_wd_fp
                  )
              }
          }
          {
            \fp_set:Nn \l__hit_quan_sx_fp
              {
                min
                  (
                    \l__hit_quan_sy_fp ,
                    \l__hit_quan_inner_fp
                      / sqrt ( \l__hit_quan_wd_fp ^ 2 + \l__hit_quan_hd_fp ^ 2 )
                  )
              }
            \fp_set_eq:NN \l__hit_quan_sy_fp \l__hit_quan_sx_fp
          }
        \fp_zero:N \l__hit_quan_dx_fp
        \fp_zero:N \l__hit_quan_dy_fp
      }
  }
%    \end{macrocode}
%
% 查表。有整组数据（一位、11、两位）就按表项放；组缺了、或者数在 0--99 以外，退到同位数的
% 那一组只取数字高，等比缩放；连那一组也没有才走通用近似。宽度系数按字号在 10bp 与 24bp
% 之间线性插值，圈线比原先粗时再乘一个余量系数。
%    \begin{macrocode}
\cs_new_protected:Npn \__hit_quan_fit:
  {
    \bool_set_false:N \l__hit_quan_fitted_bool
    \hbox_set:Nn \l__hit_quan_digits_box
      { \__hit_quan_digit_font: \str_use:N \l__hit_quan_digits_str }
    \fp_set:Nn \l__hit_quan_wd_fp { \dim_to_fp:n { \box_wd:N \l__hit_quan_digits_box } }
    \fp_set:Nn \l__hit_quan_hd_fp
      { \dim_to_fp:n { \box_ht:N \l__hit_quan_digits_box + \box_dp:N \l__hit_quan_digits_box } }
    \__hit_quan_group:
    \int_compare:nNnT { \l__hit_quan_value_int } < { 100 }
      {
        \tl_set:Nx \l__hit_quan_key_tl
          {
            \g__hit_quan_family_tl /
            \bool_if:NTF \l__hit_quan_bold_bool { bold } { regular } /
            \l__hit_quan_group_tl
          }
        \prop_get:NVNT \c__hit_quan_fit_prop \l__hit_quan_key_tl \l__hit_quan_entry_tl
          {
            \bool_set_true:N \l__hit_quan_fitted_bool
            \fp_set:Nn \l__hit_quan_tmp_fp
              {
                min ( 1 , max ( 0 , ( \l__hit_quan_fs_fp * 72 / 72.27 - 10 ) / 14 ) )
              }
            \fp_set:Nn \l__hit_quan_sx_fp
              {
                \__hit_quan_entry:n { 1 } *
                ( \__hit_quan_entry:n { 5 }
                  + ( \__hit_quan_entry:n { 6 } - \__hit_quan_entry:n { 5 } )
                    * \l__hit_quan_tmp_fp )
                * \l__hit_quan_clear_fp
              }
            \fp_set:Nn \l__hit_quan_sy_fp
              { \__hit_quan_entry:n { 2 } * \l__hit_quan_h_fp / \l__hit_quan_hd_fp }
            \fp_set:Nn \l__hit_quan_dx_fp { \__hit_quan_entry:n { 3 } * \l__hit_quan_h_fp }
            \fp_set:Nn \l__hit_quan_dy_fp { \__hit_quan_entry:n { 4 } * \l__hit_quan_h_fp }
          }
      }
    \bool_if:NF \l__hit_quan_fitted_bool
      {
        \tl_set:Nx \l__hit_quan_key_tl
          {
            \g__hit_quan_family_tl /
            \bool_if:NTF \l__hit_quan_bold_bool { bold } { regular } /
            \int_compare:nNnTF { \l__hit_quan_value_int } < { 10 } { single } { two }
          }
        \prop_get:NVNT \c__hit_quan_fit_prop \l__hit_quan_key_tl \l__hit_quan_entry_tl
          {
            \bool_set_true:N \l__hit_quan_fitted_bool
            \fp_set:Nn \l__hit_quan_tmp_fp
              { \__hit_quan_entry:n { 2 } * \l__hit_quan_h_fp / \l__hit_quan_hd_fp }
            \fp_set:Nn \l__hit_quan_sx_fp
              {
                \l__hit_quan_tmp_fp
                * min
                    (
                      1 ,
                      \l__hit_quan_avail_fp
                        / ( \l__hit_quan_tmp_fp
                            * sqrt ( \l__hit_quan_wd_fp ^ 2 + \l__hit_quan_hd_fp ^ 2 ) )
                    )
              }
            \fp_set_eq:NN \l__hit_quan_sy_fp \l__hit_quan_sx_fp
            \fp_zero:N \l__hit_quan_dx_fp
            \fp_zero:N \l__hit_quan_dy_fp
          }
      }
    \bool_if:NF \l__hit_quan_fitted_bool { \__hit_quan_fit_approx: }
  }
%    \end{macrocode}
%
% 几何。圈径、圈线粗细、内圈大小与可用宽度都以参照字形的墨迹高 $H$ 为基准：圈外缘与参照字形等大，
% 圈线中心线的直径是 $H$ 减圈线粗细；版面宽度是 $H$ 减常数圈线粗细（与圈线实际粗细无关，
% 换字体时前后字符的间距不变），两侧各留 \texttt{0.05em} 与前后字符隔开。
%    \begin{macrocode}
\cs_new_protected:Npn \__hit_quan_geometry:
  {
    \fp_set:Nn \l__hit_quan_fs_fp { \dim_to_fp:n { \f@size pt } }
    \fp_set:Nn \l__hit_quan_h_fp
      { \dim_to_fp:n { \g__hit_quan_ref_ht_dim + \g__hit_quan_ref_dp_dim } }
    \fp_set:Nn \l__hit_quan_leg_fp
      {
        0.0315 * \l__hit_quan_fs_fp * \bool_if:NTF \l__hit_quan_bold_bool { 1.45 } { 1 }
      }
    \bool_set_false:N \l__hit_quan_cal_bool
    \prop_get:NVNT \c__hit_quan_stroke_prop \g__hit_quan_family_tl \l__hit_quan_entry_tl
      {
        \bool_set_true:N \l__hit_quan_cal_bool
        \fp_set:Nn \l__hit_quan_t_fp
          {
            \l__hit_quan_h_fp
            * \__hit_quan_entry:n { \bool_if:NTF \l__hit_quan_bold_bool { 2 } { 1 } }
          }
      }
    \bool_if:NF \l__hit_quan_cal_bool { \fp_set_eq:NN \l__hit_quan_t_fp \l__hit_quan_leg_fp }
    \fp_set:Nn \l__hit_quan_w_fp { max ( 1 , \l__hit_quan_h_fp - \l__hit_quan_leg_fp ) }
    \fp_set:Nn \l__hit_quan_inner_fp
      {
        \l__hit_quan_h_fp - 2 * \l__hit_quan_t_fp
        - 2 * max ( 0.02 * \l__hit_quan_h_fp , 0.12 )
      }
    \bool_if:NTF \l__hit_quan_cal_bool
      {
        \fp_set:Nn \l__hit_quan_avail_fp
          {
            max
              (
                1 ,
                min
                  (
                    \l__hit_quan_w_fp - 0.1 * \l__hit_quan_fs_fp - \l__hit_quan_leg_fp ,
                    \l__hit_quan_inner_fp
                  )
              )
          }
      }
      {
        \fp_set:Nn \l__hit_quan_avail_fp
          {
            max
              (
                1 ,
                max ( 1 , \l__hit_quan_h_fp - \l__hit_quan_t_fp )
                  - 0.1 * \l__hit_quan_fs_fp - \l__hit_quan_t_fp
              )
          }
      }
    \fp_set:Nn \l__hit_quan_clear_fp { 1 }
    \bool_lazy_and:nnT
      { \l__hit_quan_cal_bool }
      { \fp_compare_p:nNn { \l__hit_quan_t_fp } > { \l__hit_quan_leg_fp } }
      {
        \fp_set:Nn \l__hit_quan_clear_fp
          {
            min
              (
                1 ,
                max
                  (
                    0 ,
                    \l__hit_quan_h_fp - 2 * \l__hit_quan_t_fp
                    - 2 * max ( 0.005 * \l__hit_quan_h_fp , 0.12 )
                  )
                  / max ( 0.01 , \l__hit_quan_h_fp - 2 * \l__hit_quan_leg_fp )
              )
          }
      }
  }
%    \end{macrocode}
%
% 排版：圈用 \pkg{pict2e} 画在宽度 $w$ 的盒子里，圆心在盒子水平中点、垂直落在参照字形的墨迹中点；
% 数字缩放后让墨迹中心落在圆心加偏移（垂直偏移向下为正）。盒子的高深取参照字形的，
% 不让圈把行撑高。
%    \begin{macrocode}
\cs_new_protected:Npn \__hit_quan_typeset:
  {
    \box_scale:Nnn \l__hit_quan_digits_box { \l__hit_quan_sx_fp } { \l__hit_quan_sy_fp }
    \tl_set:Nx \l__hit_quan_tmpx_tl { \fp_to_decimal:n { \l__hit_quan_w_fp / 2 } }
    \tl_set:Nx \l__hit_quan_tmpy_tl
      { \fp_to_decimal:n { \dim_to_fp:n { \g__hit_quan_ref_ht_dim - \g__hit_quan_ref_dp_dim } / 2 } }
    \tl_set:Nx \l__hit_quan_tmpd_tl
      { \fp_to_decimal:n { max ( 1 , \l__hit_quan_h_fp - \l__hit_quan_t_fp ) } }
    \hbox_set:Nn \l__hit_quan_ring_box
      {
        \unitlength = 1pt
        \linethickness { \fp_to_dim:n { \l__hit_quan_t_fp } }
        \begin{picture} (0,0)
          \put ( \l__hit_quan_tmpx_tl , \l__hit_quan_tmpy_tl )
            { \circle { \l__hit_quan_tmpd_tl } }
        \end{picture}
      }
    \tex_kern:D \fp_to_dim:n { 0.05 * \l__hit_quan_fs_fp } \scan_stop:
    \hbox_to_wd:nn { \fp_to_dim:n { \l__hit_quan_w_fp } }
      {
        \tex_vrule:D width 0pt height \g__hit_quan_ref_ht_dim depth \g__hit_quan_ref_dp_dim \scan_stop:
        \box_use:N \l__hit_quan_ring_box
        \tex_kern:D
          \fp_to_dim:n
            {
              \l__hit_quan_w_fp / 2 + \l__hit_quan_dx_fp
              - \dim_to_fp:n { \box_wd:N \l__hit_quan_digits_box } / 2
            }
        \scan_stop:
        \box_move_up:nn
          {
            \fp_to_dim:n
              {
                \dim_to_fp:n { \g__hit_quan_ref_ht_dim - \g__hit_quan_ref_dp_dim } / 2
                - \l__hit_quan_dy_fp
                - \dim_to_fp:n { \box_ht:N \l__hit_quan_digits_box - \box_dp:N \l__hit_quan_digits_box } / 2
              }
          }
          { \box_use:N \l__hit_quan_digits_box }
        \hss
      }
    \tex_kern:D \fp_to_dim:n { 0.05 * \l__hit_quan_fs_fp } \scan_stop:
  }
%    \end{macrocode}
%
% 当前宋体有现成的带圈字符就直接用，字形比画出来的规整；没有才画。
%    \begin{macrocode}
\NewDocumentCommand \quan { m }
  {
    \group_begin:
      \hit@engine@use@cjk@font:
      \int_set:Nn \l__hit_quan_slot_int { \__hit_quan_slot:n {#1} }
      \bool_lazy_and:nnTF
        { \int_compare_p:nNn \l__hit_quan_slot_int > { 0 } }
        { \hit@engine@if@cjk@glyph_p:n { \l__hit_quan_slot_int } }
        { \hit@engine@char:n { \l__hit_quan_slot_int } }
        { \__hit_quan_draw:n {#1} }
    \group_end:
  }
%    \end{macrocode}
%
% 画出来的圈由盒子和线条组成，\pkg{docxlike} 在 LuaTeX 下抬高上标靠的是字形节点的属性，
% 只抬字形、不抬线条，脚注号里的数字会和圈错开。所以把抬高量先取走、清掉属性，
% 整个圈连数字一起按盒子抬。
%    \begin{macrocode}
\int_new:N \l__hit_quan_raise_int
\box_new:N \l__hit_quan_out_box
\cs_new_protected:Npn \__hit_quan_raise_take:
  {
    \int_zero:N \l__hit_quan_raise_int
    \sys_if_engine_luatex:T
      {
        \cs_if_exist:NT \docxlike@va@attr
          {
            \int_set:Nn \l__hit_quan_raise_int { \docxlike@va@attr }
            \int_compare:nNnT { \l__hit_quan_raise_int } < { -1073741824 }
              { \int_zero:N \l__hit_quan_raise_int }
            \unsetattribute \docxlike@va@attr
            \unsetattribute \docxlike@vaja@attr
          }
      }
  }
\cs_new_protected:Npn \__hit_quan_draw:n #1
  {
    \__hit_quan_raise_take:
    \int_set:Nn \l__hit_quan_value_int { \int_abs:n {#1} }
    \str_set:Nx \l__hit_quan_digits_str { \int_use:N \l__hit_quan_value_int }
    \__hit_quan_font_ensure:
    \bool_set:Nn \l__hit_quan_bold_bool { \tl_if_head_eq_charcode_p:fN { \f@series } b }
    \__hit_quan_measure_ref:
    \__hit_quan_geometry:
    \__hit_quan_fit:
    \hbox_set:Nn \l__hit_quan_out_box { \__hit_quan_typeset: }
    \box_move_up:nn { \l__hit_quan_raise_int sp } { \box_use:N \l__hit_quan_out_box }
  }
%    \end{macrocode}
% \end{macro}
%

% \begin{macro}{\hit@setup@footnote}
%    \begin{macrocode}
\cs_new_protected:Npn \hit@setup@footnote
  {
%    \end{macrocode}
%
% 字体缺少带圈数字时用 \pkg{pict2e} 画圈，画法见上面的 \cs{quan}。
%    \begin{macrocode}
\RequirePackage{pict2e}
%    \end{macrocode}
%
% 脚注编号与脚注样式。
%    \begin{macrocode}
\cs_new_protected:Npn \hit@text@circled ##1 { \quan { \arabic {##1} } }
\cs_set:Npn \thefootnote   { \hit@text@circled { footnote } }
\cs_set:Npn \thempfootnote { \hit@text@circled { mpfootnote } }
\cs_set:Npn \footnoterule
  { \vskip -3\p@ \hrule \@width 0.3\textwidth \@height 0.4\p@ \vskip 2.6\p@ }
\cs_set_eq:NN \hit@original@footnote@size \footnotesize
\cs_set:Npn \footnotesize { \hit@original@footnote@size \hit@fontsize@x{xiaowu}{1.5} }
\cs_if_exist:NT \footnotemargin { \footnotemargin1.5em\relax }
%    \end{macrocode}
%
% 脚注正文中的标号按正文大小排版，不使用上标。
%    \begin{macrocode}
\cs_set_eq:NN \hit@original@make@footnote@mark \@makefnmark
\cs_new:Npn \hit@footnote@mark { \hbox { { \normalfont \@thefnmark } } }
\cs_set_eq:NN \hit@original@make@footnote@text \@makefntext
\cs_set:Npn \@makefntext ##1
  {
    \cs_set_eq:NN \@makefnmark \hit@footnote@mark
    \hit@original@make@footnote@text {##1}
    \cs_set_eq:NN \@makefnmark \hit@original@make@footnote@mark
  }
  }
%    \end{macrocode}
% \end{macro}
%
%    \begin{macrocode}
%</shared-footnote>
%    \end{macrocode}
%
% \Finale
